% TOP_PROBLEM: {"id": 30007092, "title": "Power improvement over exponent two in Heilbronn triangle constructions", "category_id": 2, "category": {"id": 2, "name": "combinatorics", "display_name": "Combinatorics", "description": "Counting problems, graph theory, discrete structures.", "slug": "combinatorics", "order_index": 2, "created_at": "2026-07-31T15:26:25.671Z"}, "status": "open"}
% FIRST_POSTED: 2026-10-06
% ATTEMPT_STATUS: solved
% OPENAI_CLAIM: {"schema_version":1,"source_repo":"https://github.com/openai/math","source_commit":"adc7f1241b42e322a6451854ab7e4b4c146bf78a","release_date":"2026-10-06","match":"full","resolution":"proved","independently_reviewed":false,"second_pass":"agreed","families":[{"family":"191","title":"A power improvement in the Heilbronn triangle lower bound","subject":"Combinatorics","manuscripts":[{"title":"A power improvement in the Heilbronn triangle lower bound","pdf_path":"preprints/A-power-improvement-in-the-Heilbronn-triangle-lower-bound-September-25-2026/main.pdf","readme_path":"preprints/A-power-improvement-in-the-Heilbronn-triangle-lower-bound-September-25-2026/README.md","date":"2026-09-25","theorem_ref":"intro:main (sections/01-introduction.tex:42)","statement_sources":[{"path":"preprints/A-power-improvement-in-the-Heilbronn-triangle-lower-bound-September-25-2026/build/sections/01-introduction.tex","line":42,"label":"intro:main"}]}],"lean":{"doc_path":"lean/docs/191.md","comparators":["lean/ComparatorChallenges/HeilbronnTriangle.json","lean/ComparatorChallenges/HeilbronnTriangle.lean"],"declarations":[],"covers_main_theorem":false},"reasoning_trace":null}]}
\documentclass[11pt]{article}
\usepackage[margin=25mm]{geometry}
\usepackage{fontspec}
\setmainfont{Latin Modern Roman}
\usepackage{amsmath,amssymb,mathtools}
\usepackage{unicode-math}
\setmathfont{Latin Modern Math}
\usepackage{enumitem}
\usepackage{xurl,hyperref,bookmark}
\hypersetup{colorlinks=true,urlcolor=blue}
\setcounter{secnumdepth}{0}
\setlength{\parindent}{0pt}
\setlength{\parskip}{5pt}
\setlength{\emergencystretch}{3em}
\begin{document}
\section*{Power improvement over exponent two in Heilbronn triangle constructions}

\subsection{OpenAI's claim}
\textit{Heilbronn triangle power improvement}: For every sufficiently large $n$, the paper claims $n$ points in the unit square with every triangle area at least $c n^{-2+\eta}$ for fixed positive $c,\eta$. This contradicts an upper bound with exponent arbitrarily close to minus two.

\subsection{Mathematical statement of the claimed result}
\paragraph{Mathematical claim 1.} For $n\ge3$, define $\Delta(n)=\max_{P\subset[0,1]^2,\ |P|=n}\min_{\{p,q,r\}\subset P}\tfrac12|\det(q-p,r-p)|$. The claim gives absolute $\eta,c>0,n_0$ with \[\Delta(n)\ge cn^{-2+\eta}\quad(n\ge n_0).\] It therefore refutes the proposed upper estimate $\Delta(n)=O_\epsilon(n^{-2+\epsilon})$ for every $\epsilon>0$, by taking $\epsilon<\eta$. A fixed affine scaling transfers the construction to a unit disk.

\subsection{Outline of the argument}
\textit{Heilbronn triangle power improvement}: Two congruence conditions govern randomly sampled integer columns. Norm-polynomial residues control small nonzero determinants, while an independent quadratic-surface condition and lattice counts handle zero determinants. The Chinese remainder construction combines them, repeated labels supply a deletion saving, and scaling transfers the retained determinant lower bound to triangle areas.

\subsection{Scope relative to this problem}
\textit{Heilbronn triangle power improvement}: The manuscript’s result states Delta(n)>=c\_1 n\textasciicircum{}(-2+eta) for every sufficiently large n, with Delta the supremal smallest triangle area in the unit square. Taking a slightly smaller c supplies actual point-set witnesses even if the supremum convention is used.

With the hypotheses and conventions above, the claimed result addresses the whole mathematical assertion on this page.

\subsection{Formal verification}
Family 191 describes Lean coverage for the triangle-area improvement along an unbounded sequence of sizes; the every-large-size extension is broader.

\textit{Heilbronn triangle power improvement}: The manuscript links \texttt{HeilbronnTriangle.lean}.

The documented selection does not cover the entire manuscript assertion relevant here.

This site has not run the Lean development. These links describe OpenAI’s selected formal statements and should be read with their stated scope.

\subsection{Status caveat}
OpenAI reports that an internal model produced the manuscripts. Their arguments have not been independently verified as proofs by this site. Any portion outside the documented formal scope may contain mistakes, and the catalogue label reports the announced mathematical claim.

\subsection{Sources}
\begin{itemize}
\item OpenAI, \emph{A power improvement in the Heilbronn triangle lower bound}, 2026-09-25; citation key \texttt{OAI:A-power-improvement-in-the-Heilbronn-triangle-lower-bound-September-25-2026}. \href{https://github.com/openai/math/blob/adc7f1241b42e322a6451854ab7e4b4c146bf78a/preprints/A-power-improvement-in-the-Heilbronn-triangle-lower-bound-September-25-2026/main.pdf}{PDF (GitHub)}; \href{https://github.com/openai/math/raw/adc7f1241b42e322a6451854ab7e4b4c146bf78a/preprints/A-power-improvement-in-the-Heilbronn-triangle-lower-bound-September-25-2026/main.pdf}{PDF (direct)}; \href{https://github.com/openai/math/blob/adc7f1241b42e322a6451854ab7e4b4c146bf78a/preprints/A-power-improvement-in-the-Heilbronn-triangle-lower-bound-September-25-2026/README.md}{Manuscript page and citation}; \href{https://github.com/openai/math/blob/main/preprints/A-power-improvement-in-the-Heilbronn-triangle-lower-bound-September-25-2026/main.pdf}{Latest version}; Statement source: \href{https://github.com/openai/math/blob/adc7f1241b42e322a6451854ab7e4b4c146bf78a/preprints/A-power-improvement-in-the-Heilbronn-triangle-lower-bound-September-25-2026/build/sections/01-introduction.tex\#L42}{intro:main}.
\item Family 191: \href{https://github.com/openai/math/blob/adc7f1241b42e322a6451854ab7e4b4c146bf78a/lean/docs/191.md}{Lean scope}.
\item \href{https://github.com/openai/math/blob/adc7f1241b42e322a6451854ab7e4b4c146bf78a/lean/ComparatorChallenges/HeilbronnTriangle.json}{Comparator HeilbronnTriangle.json}.
\item \href{https://github.com/openai/math/blob/adc7f1241b42e322a6451854ab7e4b4c146bf78a/lean/ComparatorChallenges/HeilbronnTriangle.lean}{Comparator HeilbronnTriangle.lean}.
\item \href{https://github.com/openai/math/blob/adc7f1241b42e322a6451854ab7e4b4c146bf78a/CONTENTS.md}{OpenAI catalogue entry}.
\end{itemize}
\end{document}
